% =====================================================================
% COSC 3337 — Homework 4 — ANSWER TEMPLATE for Part 1
%
% Typesetting Part 1 in LaTeX earns 5 bonus points.
%
% To use this:  upload it to Overleaf (or run pdflatex locally), put your name
% in the \student command, fill in the answers, and submit the PDF.
% You do not submit this .tex file.
%
% ---------------------------------------------------------------------
% WRITING CODE AND SQL
%
% Put code (and anything else you want shown exactly, character for character)
% in a code block:
%
%   \begin{lstlisting}
%   SELECT * FROM ticket_sale_wide LIMIT 5
%   \end{lstlisting}
%
% Inside a code block, everything is printed exactly as you type it: no
% escaping needed.
%
% For a short piece of code inside a sentence, use \lstinline|...| :
%
%   The query \lstinline|SELECT COUNT(*) FROM ticket_sale_wide| counts rows.
%
% Outside code, underscores must be written \_ , as in ticket\_sale\_wide.
%
% Math goes between dollar signs, e.g.
%   $\{A, B\}^+ = \{A, B, C\}$     $\alpha \rightarrow \beta$     $X \oplus Y$
%   $R_1 \cap R_2$     and \underline{A} to underline an attribute
% =====================================================================
\documentclass[11pt]{article}

\usepackage[margin=1in]{geometry}
\usepackage{amsmath,amssymb}
\usepackage{booktabs}
\usepackage{enumitem}
\usepackage{fancyhdr}
\usepackage{textcomp}
\usepackage{listings}

\lstset{
  basicstyle=\ttfamily\small,
  columns=fullflexible,
  keepspaces=true,
  upquote=true,
  showstringspaces=false,
  breaklines=true,
  frame=leftline,
  xleftmargin=1.2em,
}

\newcommand{\student}{YOUR NAME HERE}

\pagestyle{fancy}
\fancyhf{}
\lhead{\small COSC 3337 — HW 4 — \student}
\rhead{\small Page \thepage}
\renewcommand{\headrulewidth}{0.4pt}

\setlength{\parindent}{0pt}
\setlength{\parskip}{6pt}

% =====================================================================
\begin{document}

\begin{center}
  {\Large\bfseries COSC 3337 — Homework 4}\\[2pt]
  {\large \student}
\end{center}


% =====================================================================
\section*{1.1 RAID}

\textbf{(a)}

\begin{tabular}{@{}lll@{}}
\toprule
Layout & Usable capacity & Failures always survived \\
\midrule
RAID 0 & ... & ... \\
RAID 1 & ... & ... \\
RAID 5 & ... & ... \\
\bottomrule
\end{tabular}

\textbf{(b)}
\begin{enumerate}
  \item
\begin{lstlisting}
your work here
\end{lstlisting}
  \item
\begin{lstlisting}
your work here
\end{lstlisting}
\end{enumerate}

\textbf{(c)}

% =====================================================================
\section*{1.2 Slotted pages}

\begin{tabular}{@{}lccc@{}}
\toprule
Slot & (offset, size) now & After (a) & After (b) \\
\midrule
1 & (215, 25) & ... & ... \\
2 & (175, 40) & ... & ... \\
3 & (160, 15) & ... & ... \\
4 & (130, 30) & ... & ... \\
5 & ---       & ... & ... \\
\midrule
End of free space & 130 & ... & ... \\
Free bytes        & 110 & ... & ... \\
\bottomrule
\end{tabular}

Work for (a):

Work for (b):

% =====================================================================
\section*{1.3 File organization}

\textbf{(a)}

\textbf{(b)}

\textbf{(c)}

% =====================================================================
\section*{1.4 Functional dependencies in the export}

\textbf{(a)}

\textbf{(b)}

\textbf{(c)}

% =====================================================================
\section*{1.5 BCNF decomposition}

\textbf{(a)}

\textbf{(b)}

\begin{enumerate}
  \item ...
\end{enumerate}

Final relations:

\textbf{(c)}

\end{document}
